\documentclass[sigconf,nonacm]{acmart}
\settopmatter{printacmref=false}
\setcopyright{none}
\renewcommand\footnotetextcopyrightpermission[1]{}
\pagestyle{plain}
\raggedbottom

\acmConference[Manuscript Draft]{Layer-Safe Generative Design Study}{2026}{}
\acmYear{2026}
\acmISBN{}
\acmDOI{}

\title{Layer-Safe Generative Design from Nwagu Aneke: A Non-Promotion Invariant}
\author{Anonymous Author(s)}
\affiliation{\institution{Independent research manuscript}\country{Nigeria}}
\email{anonymous@example.invalid}

\begin{abstract}
Artifact-derived design systems can generate useful interfaces, grammars, and
software abstractions from cultural source material. They can also create a
serious epistemic failure: generated outputs may promote derived, speculative,
or blocked interpretations into source-observed claims. This article formulates
and tests a non-promotion invariant for design systems derived from the Nwagu
Aneke evidence package. The accepted source layer is 26 rows by 8 vowel or
modifier columns, yielding 208 records; 27/216 is a derived f/v split layer
only. The experiment evaluates 24 layer-safety cases, rejects all 14
intentional unsafe promotions, and compares the invariant with provenance-only,
citation-only, and no-label baselines. A ten-case transfer and counterexample
audit then tests the same rule across digital-edition, annotation,
Unicode-process, cultural-heritage machine-learning, and rights-boundary
examples. The accompanying finite-composition lemma states that a pipeline
composed only of layer-safe transforms cannot upgrade a non-source claim into a
source-observed claim. The contribution is a specific systems result for
evidence-preserving generative design; universal compression, glyph-shape
completion, and release authorization remain outside its scope.
\end{abstract}

\keywords{Nwagu Aneke, generative design, evidence layers, claim safety, provenance, cultural heritage systems}

\begin{document}
\maketitle

\section{Introduction}
The Nwagu Aneke artifact has a stable working foundation for the present study:
source-observed 26 rows by 8 vowel or modifier columns equals 208 records, while
27/216 is a derived f/v split layer only. That foundation creates a design
opportunity. It also creates a design risk. Once an artifact enters a generative
workflow, outputs can look authoritative even when their evidence layer is weak.
A design grammar may turn a derived normalization into an apparent source fact.
A visualization may hide a blocked assumption. A paper draft may polish an
unsupported claim until it reads like a result.

This article asks whether that failure can be prevented by design. The
research question is: can artifact-derived generative systems enforce a non-promotion
invariant, so that a generated output never claims a stronger evidence layer
than its inputs license? The paper tests the question using the Nwagu Aneke
count-layer distinction as the motivating case. It does not ask the reader to
accept a complete PAGC theory. It asks whether a finite set of transform rules
can block a specific class of epistemic error.

Hypothesis: if evidence layers are treated as ordered claim types and every
generative transform is required to preserve or weaken, rather than strengthen,
the input layer without a review event, then unsafe promotions from derived,
speculative, or blocked claims into source-observed claims can be detected while
safe source-preserving and derived-preserving outputs remain usable. This
hypothesis is deliberately narrower than a general trust-in-AI claim. It is a
testable systems claim about artifact-derived design pipelines.

The answer is positive within the tested scope. An expanded 24-case
layer-safety test identified and rejected all 14 intentional unsafe promotions.
Provenance-only, citation-only, and no-label baselines missed promotion errors
that the layer rule caught. A finite composition lemma then states the general
systems property: if every transform
in a finite pipeline is layer-safe, their composition is layer-safe. This is a
small theorem, but it matters because it converts a recurring workflow failure into a
testable design invariant.

The invariant also connects this cultural-heritage problem to adjacent
representational systems. Unicode proposal guidance asks for evidence about
characters and usage, while tokenization research shows that representational
units affect downstream model behavior
\cite{seiProposalTips,sennrich2016bpe,kudo2018sentencepiece,bostrom2020bytepair,rust2021tokenization}.
This paper does not solve encoding or tokenization. It supplies a prior rule:
before a generated interface, tokenizer, or article draft can use an artifact
claim, the system must know whether that claim is source-observed, derived,
speculative, or blocked.

\section{Related Work}
The domain anchor is Azuonye's account of the Nwagu Aneke script, supplemented
by public script-status records and source catalogues \cite{azuonye1992,omniglotNwagu,unicodeAfricanScripts2023,scriptSourceNwagu,sirisNwaguProposal}.
The systems anchor comes from provenance, research-object packaging, annotation,
and cultural-heritage representation: PROV-O, RO-Crate, TEI, IIIF, Web
Annotation, CIDOC CRM, DataCite, FAIR, and CARE \cite{w3cProvO,rocrate12,teiP5Glyphs,iiifPresentation3,w3cAnnotationModel,w3cAnnotationVocab,cidocCrm,datacite45,wilkinson2016fair,carroll2020care}.
The evaluation anchor comes from scientific claim verification, citation-grounded
generation, and autoresearch loops \cite{wadden2020scifact,gao2023alce,karpathyAutoresearch}.

These bodies of work make the present result more modest and more precise.
Provenance standards can record derivation, but recording derivation is not the
same as preventing a generated output from overclaiming it. Citation systems can
require sources, but a source can be cited for the wrong layer. Research-object
packaging can preserve files, but it does not decide whether a design layer has
been promoted into a source layer. The contribution here is therefore not a new
metadata standard. It is a non-promotion invariant for artifact-derived
generative workflows.

\section{Materials and Evidence}
The evidence package is the same bounded Nwagu Aneke foundation used by the
count-layer audit: source-observed 26 by 8 equals 208; derived 27/216 only by
f/v split. The experiment does not use public source-image reproduction and does
not depend on a completed glyph-shape grammar. It tests layer behavior over
claims and transforms.

The input layer set is ordered by permissible assertion strength. A
source-observed claim can support direct source statements. A derived claim can
support analytical or design statements when labeled. A speculative claim can
support hypotheses and experiment proposals. A blocked claim can support only
limitations, requirements, or future work. A transform is unsafe if it emits a
claim at a stronger layer than the input evidence permits.


\begin{table}[t]
\caption{Evidence and claim-status summary}
\label{tab:layer_safety}
\begin{tabular}{p{0.34\columnwidth}p{0.56\columnwidth}}
\toprule
Item & Status \\
\midrule
Source foundation & 26 rows by 8 columns equals 208 records \\
Derived foundation & 27/216 by f/v split only \\
Experiment & 24 layer-safety cases with shallow-baseline comparison \\
Unsafe promotions & 14 intentional promotions rejected \\
Maximum claim & Systems invariant, not glyph grammar or universal theory \\
\bottomrule
\end{tabular}
\end{table}


\section{Method}
The method defines a layer partial order and a transform rule. Let layers be
ordered from strongest to weakest as source-observed, derived, speculative, and
blocked for public-claim purposes. A transform is layer-safe when its output
layer is no stronger than the strongest layer licensed by its inputs. In plain
terms, a function may preserve or weaken certainty, but it may not strengthen it
without an explicit review event.

The test set contains safe and unsafe examples. Safe examples preserve the
source layer, describe 27/216 as derived, safely weaken claim status, or move
weak claims into limitations. Unsafe examples describe the derived matrix as the
primary source layer, treat design grammar as historical evidence, or turn an
unresolved source question into a publication claim. The evaluation asks whether
the non-promotion rule rejects the unsafe cases while allowing useful derived
design work, then compares the rule with provenance-only, citation-only, and
no-label baselines.

The proof artifact is finite and operational. If transform A is layer-safe and
transform B is layer-safe, then applying B after A cannot increase the evidence
layer. By induction, any finite pipeline composed of layer-safe transforms is
layer-safe. This is not a claim that the generated output is true. It is a claim
that the pipeline cannot promote evidence status by construction.

\section{Results}
The expanded 24-case experiment rejected all 14 intentional unsafe promotions.
The non-promotion rule allowed source-preserving, derived-label-preserving, and
safe-weakening outputs, while blocking outputs that upgraded derived,
speculative, or blocked claims into stronger evidence layers. The shallow
baselines missed the target failure mode: provenance-only checking allowed 11
promotion errors, citation-only checking allowed 7, and a no-label filter allowed
all 14. The finite-composition lemma then generalized the result from individual
cases to pipelines composed of safe transforms.

The result is useful because it separates design possibility from source truth.
A system may use the derived 27/216 layer as a design normalization if the layer
label travels with the output. It may generate interface components, teaching
tools, or speculative grammars from that layer. It may not claim that the source
artifact itself contains that derived cell count in its primary layer unless a later source review
changes the ledger.

\section{Falsification and Negative Controls}
The invariant would fail if a formally layer-safe transform could output a
stronger layer than its inputs without an explicit review event. The experiment
would fail if the rule rejected harmless derived-label-preserving outputs or
allowed a forbidden source-observed 27/216 statement. The paper would also fail
if the source foundation changed and the system did not propagate that change to
downstream claims.

The negative controls are important. A manually written design note without
layer labels is not enough. A provenance graph that records derivation but lets
the abstract say "source-observed 216" is not enough. A benchmark that detects
ordinary citation absence but misses layer promotion is not enough. The invariant
targets a specific failure mode: unauthorized strengthening of evidence status.

\section{Discussion}
The systems contribution is small but reusable. Many research workflows already
use provenance, citations, metadata, and review. The missing control is often a
claim-type discipline that survives generation. Nwagu Aneke makes the problem
visible because the 26 by 8 and 27/216 layers are easy to confuse and attractive
to formalize. The same structure appears in other artifact-derived workflows:
edition versus interpretation, image crop versus character, interface token
versus source sign, public claim versus blocked evidence.

For design, the implication is constructive. The invariant does not prevent
experimentation. It enables safer experimentation. A generative system can
explore derived grammars, layouts, interface states, and teaching keys, provided
that every output carries its layer and maximum public claim. The invariant lets
researchers investigate applications from the accepted foundation without
turning every application into a source claim.


\section{Relation to Provenance and Type Systems}
The non-promotion invariant is related to provenance but not identical to it.
PROV-O can describe that an output was generated by an activity from an input
entity \cite{w3cProvO}. RO-Crate can package that activity and its files
\cite{rocrate12}. Web Annotation can attach an interpretive body to a
source target \cite{w3cAnnotationModel,w3cAnnotationVocab}. These are
necessary capabilities, but they do not by themselves prevent a generated
sentence from saying more than its evidence licenses. The invariant adds a
claim-strength rule on top of provenance.

The invariant is closer to a type discipline. A claim has an evidence type:
primary layer, derived layer, speculative layer, or blocked layer. A transform
has a type signature that constrains what it may emit. If a transform receives a
derived claim, it may emit a derived design output or a weaker limitation, but
it may not emit a direct primary-layer claim. This type-like framing makes the
system inspectable. A reviewer can ask where the layer was introduced and
whether any function strengthened it without review. The broader standards
context includes script documentation, source presentation, annotation,
provenance, research-object packaging, data citation, open-science reuse, and
community-governance constraints \cite{azuonye1992,omniglotNwagu,unicodeAfricanScripts2023,scriptSourceNwagu,sirisNwaguProposal,teiP5Glyphs,iiifPresentation3,w3cAnnotationModel,w3cAnnotationVocab,w3cProvO,rocrate12,cidocCrm,datacite45,wilkinson2016fair,carroll2020care,wadden2020scifact,gao2023alce,karpathyAutoresearch,acmSubmissions,arxivSubmitTex}.

\section{Design Consequences}
Layer safety changes how artifact-derived design systems are built. A visual
atlas should not present all nodes with the same epistemic weight. A prompt
template should not ask a model to prove applications when the evidence only
supports hypotheses. A product interface should not show a derived grammar as if
it were the original script. A paper generator should not let an abstract make a
stronger claim than the results section. These are design failures, not only
writing failures.

The Nwagu Aneke foundation is useful because it gives the system a concrete
stress test. A safe system can say that 26 by 8 equals 208 in the primary
layer. It can say that 27/216 is derived through f/v splitting. It can propose
interfaces, grammars, and experiments from the derived layer. It cannot erase
the word "derived" when the output becomes public-facing. The invariant is
therefore a practical rule for preserving research integrity through generation,
formatting, visualization, and publication.

\section{Comparison to Claim Verification}
Scientific claim-verification work asks whether a claim is supported or refuted
by evidence \cite{wadden2020scifact,gao2023alce}. Layer promotion is a
related but distinct problem. A sentence may be loosely supported by a source
while still making the wrong kind of claim. For example, a source may support
that a chart exists, while a derived table supports a normalized design count.
If a model cites the chart for the normalized count, ordinary citation presence
is insufficient.

The contribution of the invariant is to require the support relation to preserve
evidence type. A source can support a primary-layer claim. A derived
transformation can support a derived-layer claim. A speculative analogy can
support a hypothesis. A blocked source question can support a limitation. The
system fails when these support relations are crossed.

\section{Use in Autonomous Research}
Autoresearch loops are useful when they propose an intervention, run a bounded
experiment, score the result, and keep or reject the change \cite{karpathyAutoresearch}.
In a cultural-artifact workflow, the score cannot be only model accuracy or
formatting quality. It must include evidence preservation. A generated paper
that looks polished but upgrades weak claims has lower research quality, not
higher.

The non-promotion invariant supplies one metric for such loops. A revision is
kept only if it preserves or weakens claim strength unless a review event
authorizes promotion. A manuscript generator, benchmark generator, visual atlas,
or interface builder can be tested on this rule. The rule is not sufficient for
truth, but it is necessary for disciplined research automation.

\section{Ethical and Community Boundary}
The invariant also protects community and rights boundaries. A blocked claim
about public circulation should remain blocked until authority and rights review
passes. A design system should not turn a private or sensitive source artifact
into a public dataset by changing format. A model should not infer permission
from access. CARE and FAIR principles therefore operate together here:
reusability is valuable, but cultural authority and responsibility constrain
what may be reused \cite{wilkinson2016fair,carroll2020care}.

For Nwagu Aneke, the present evidence supports textual claims about layer
behavior. It does not authorize source-image circulation, manuscript-publication
claims, or community-facing standardization. That boundary does not weaken the
systems result. It demonstrates the invariant: blocked circulation claims remain
blocked even when a technical rule is reproducible.



\section{Scope Conditions}
The non-promotion invariant applies under explicit scope conditions. First, the artifact must be
represented as a set of typed evidence layers rather than as a single flattened
table. Second, every transformation must preserve a record of which layer it
uses. Third, the publication workflow must treat layer labels as part of the
claim, not as optional metadata. Fourth, the rights and authority boundary must
remain active even when the technical result is reproducible.

These conditions are demanding, but they are realistic. Many cultural-heritage
and digital-humanities workflows already separate source, transcription,
interpretation, and edition. What this paper adds is the insistence that
downstream computational and publication systems preserve those separations. A
chart, table, database, prompt, interface, or manuscript abstract must not erase
the layer distinction simply because the derived representation is cleaner.

The Nwagu Aneke case is therefore a useful test object. It has a source-layer
count that can be stated, a derived f/v split that can also be useful, and a
history of speculative downstream claims that become risky when the layer
boundary is not enforced. That combination makes the result relevant beyond the
single artifact. The general lesson is that artifact-derived systems need
epistemic typing before they need more polished outputs.

\section{Alternative Explanations}
One alternative explanation is that the apparent conflict is merely a notation
problem. On that view, one could write a footnote explaining that 26 by 8 and
27/216 are different conventions and proceed normally. This paper rejects that
solution because the archive shows that notation drift can become claim drift.
Once a derived convention enters a title, abstract, figure, table, or benchmark,
readers may no longer see the footnote. The layer distinction must therefore be
represented structurally.

Another alternative explanation is that the problem is ordinary citation
failure. A better citation checker would catch unsupported statements. Citation
checking is necessary but incomplete. A claim can cite a real source and still
misstate the layer supported by that source. The issue is not only whether there
is evidence; it is whether the evidence supports the exact kind of claim being
made. Scientific claim-verification systems and citation-grounded generation
work make this problem visible, but the Nwagu Aneke case adds a cultural-artifact
layer model that those systems must preserve \cite{wadden2020scifact,gao2023alce}.

A third alternative explanation is that the derived layer should simply be
discarded. That would be too conservative. Derived representations are often
useful for design, teaching, normalization, search, and experiment construction.
The goal is not to ban derived layers. The goal is to keep them honest. A
derived layer can be productive if it remains marked as derived and if every
application states its maximum public claim.

\section{Article-Level Contribution}
The article contributes to generative design, autonomous research systems, and provenance-aware software. Its contribution is not a decorative
application of Nwagu Aneke to a fashionable field. It is a controlled result:
a design system can use artifact-derived structure without upgrading weak claims into source claims. The result matters because underdocumented scripts often become
visible to computation through partial, uncertain, or mediated records. If the
first computational representations erase uncertainty, later systems inherit a
false confidence.

The contribution also matters for AI-assisted research. Generated text is good
at producing coherent prose from partial structures. That strength becomes a
failure when prose smooths over weak evidence. A research loop that writes
without layer discipline can convert a working hypothesis into an apparent
finding. The evaluation standard used here therefore treats clarity, citation, and
format as insufficient. The paper must show which result was produced, what it
does not prove, and what would force revision.

\section{Future Experiment Design}
The next experiment should not be another broad paper-generation pass. It should
be a narrow test with a negative control. For the count-layer line of work, the
next test is an independent transcription review in which reviewers label the
row and column inventory without seeing the derived design interpretation. The
negative control is a shuffled or merged-row chart that tests whether reviewers
overfit to the expected answer. The outcome should be a disagreement matrix,
not a polished paragraph.

For the layer-safety line of work, the next test is a model comparison over
claims. Agents should receive the same evidence package under different prompt
conditions: no layer labels, explicit layer labels, claim-filter instructions, and
adversarial review instructions. The metric should not be generic fluency. It
should be promotion-error precision, promotion-error recall, severity-weighted
recall, and derived-layer preservation. The negative control is a set of claims
with trigger words but no actual layer promotion.

Those experiments would move the method toward a stronger venue submission. The
present manuscript has a coherent result, clear boundaries, and named human
review questions.

\section{Transfer and Counterexamples}
The transfer audit asks whether the invariant is only a local Nwagu count rule
or whether it also captures related cultural-heritage representation failures.
It adds ten compact cases across digital editions, annotation-region claims,
Unicode-process statements, cultural-heritage machine-learning readings, and
rights or authority blockers. Six cases intentionally upgrade weaker evidence
into source-observed claims. Four cases preserve their layer and should remain
usable.

The rule rejects all six transfer promotions and allows all four
layer-preserving cases. The counterexamples are important because several unsafe
outputs still contain a provenance link or citation. A derived editorial
normalization can cite its source and still become false if it is reported as
direct transcription. A region annotation can be traceable and still overstate
the whole object. A script-interest citation can exist while a completed
repertoire claim remains unsupported. The invariant therefore adds a
claim-strength constraint that provenance and citation alone do not supply.

\section{Evaluation Rubric}
The evaluation rubric has six criteria. Evidence strength asks whether the
central claim is supported by named artifacts rather than by prose alone.
Layer preservation asks whether every source, derived, speculative, and blocked
claim keeps its label. Reproducibility asks whether a reviewer can locate the
inputs, result file, and decision file. Citation sufficiency asks whether the
paper is situated in source-critical, standards, provenance, and claim
verification literatures. Rights clarity asks whether the paper avoids
unauthorized release claims. Reviewer survivability asks whether the paper names
the strongest reasons it could be rejected.

The rubric intentionally gives no credit for visual polish by itself. A
beautiful figure that collapses evidence layers is harmful. A formatted PDF that
turns a blocked claim into a contribution is harmful. A long bibliography that
does not support the precise claim is harmful. The paper must therefore be
readable, reproducible, and constrained by the evidence layer it inherits.

\section{Worked Example}
Consider a downstream system that wants to build a teaching interface. It may
begin with the source-layer grid and display the 26 by 8 structure as the
currently accepted source index. It may also offer a derived f/v split view for
pedagogical comparison. A safe interface labels those views differently and
keeps the transformation visible. An unsafe interface silently replaces the
source view with the derived view and then reports the derived count as the
historical inventory.

The same pattern applies to article writing. A safe abstract says that the
source layer is 26 by 8 and that the f/v expansion is derived. An unsafe abstract
reports only the cleaner derived matrix. A safe related-work section cites
standards as infrastructure and uses the artifact audit as the local result. An
unsafe related-work section treats standards citations as proof that the local
count claim is correct. A safe conclusion states what was proven and what was
not. An unsafe conclusion turns future applications into present findings.

\section{Field-Facing Reuse}
The immediate reusable output is not a public glyph dataset. It is a disciplined
pattern for handling unstable evidence layers in artifact-derived systems. A
digital-humanities project can reuse the pattern for edition layers. A
low-resource NLP project can reuse it when a culturally grounded tokenizer is
derived from incomplete source evidence. A cultural-heritage graph project can
reuse it to separate object metadata, interpretation events, and publication
claims. An autonomous-research benchmark can reuse it to score whether agents
preserve or promote claim layers.

This reuse claim is intentionally modest. The paper does not assert that every
project should adopt this paper's terminology. It asserts that the failure mode is
real enough to require a control. The local names may change, but the boundary
between observed evidence, derived representation, hypothesis, and blocked
release should not disappear.

\section{External Review Positioning}
For external review, the paper should be positioned as a bounded methods result
rather than as a comprehensive account of Nwagu Aneke. A reviewer in digital
humanities should be able to evaluate whether the source/derived distinction is
documented and whether the article avoids overclaiming. A reviewer in
information systems should be able to evaluate whether the governance mechanism
is explicit enough to prevent claim drift. A reviewer in human-computer
interaction or design should be able to evaluate whether the proposed control
changes the behavior of generated interfaces or manuscripts. A reviewer in
African writing systems should be able to see what is claimed about the artifact
and what is explicitly not claimed.

This positioning matters because the paper is interdisciplinary but not
unbounded. It does not ask every reviewer to accept every adjacent field. It
offers each field a different inspection point. The artifact scholar inspects
the source boundary. The systems reviewer inspects the invariant or audit
method. The standards reviewer inspects interoperability. The ethics reviewer
inspects rights and authority boundaries. The AI reviewer inspects whether the
workflow reduces layer-promotion errors. A paper that can be attacked along
specific lines is stronger than one whose contribution dissolves into broad
language.

\section{Failure Cases That Remain Out of Scope}
Several important failures remain outside the scope of the result. The paper does not
evaluate model performance on source images. It does not measure human reading
accuracy. It does not compare competing glyph-shape grammars. It does not
provide a public corpus. It does not settle dialect, orthographic, or
manuscript-historical questions. It does not claim that the derived layer is
better for teaching or software than the primary layer. It only states the layer
conditions under which such studies could be conducted.

Leaving these failures out of scope is not evasion. It is stage discipline. If a
future paper studies OCR, it needs image rights, crops, labels, baselines, and
error analysis. If a future paper studies tokenization, it needs corpora,
baselines, downstream tasks, and confidence intervals. If a future paper studies
Unicode readiness, it needs repertoire stability, glyph evidence, usage
examples, names, and community review. This paper provides a control condition
for those studies; it does not replace them.

\section{Scope Rationale}
The article's claim is deliberately bounded. Its evidence is named, its
limitations are active, and its central result can be checked against finite
case matrices. The manuscript avoids broad claims about universal compression,
exceptional mathematics, public source circulation, or completed
standardization. The result is a methods claim: layer-safe transforms can block
unauthorized evidence promotion while preserving useful derived design work.

The transfer audit sharpens that rationale. The invariant is not valuable
because it names a local bookkeeping preference. It is valuable because the same
failure recurs when editorial normalizations, annotation regions, OCR-derived
readings, standards interest, or rights-blocked records are allowed to speak as
though they were source-observed facts. The paper's scope is therefore broad
enough for cultural-heritage computing review and narrow enough to remain
falsifiable.

\section{Boundary Conditions for Future Submission}
Future submission work requires a
different review step. The first condition is venue-specific prior art: the paper must
be revised against the expectations of one target venue rather than remaining
generically interdisciplinary. The second condition is final source review:
claims about the artifact must be checked against the source dossier and any
available primary records. The third condition is rights and authority review:
the submission must identify exactly what source material, if any, becomes
public. The fourth condition is independent reviewer stress testing: at least one
reviewer should try to reject the paper on novelty, method, evidence, and ethics
grounds.

These conditions are the route to responsible publication. They preserve the
separation between a reproducible methods result and the social, legal, and
source-critical decisions required for external circulation.

\section{Practical Checklist for Reuse}
The practical checklist is short. A researcher reusing this paper should first
identify the layer of every claim. They should then list the transformation that
created any derived object. They should state the maximum public claim allowed
by that layer. They should add a negative control that would catch accidental
promotion. They should write limitations before writing contribution language.
They should record rights and authority assumptions separately from technical
results. Finally, they should rerun the layer-safety check after every major
manuscript revision.

This checklist is deliberately operational. It translates the argument into a
repeatable workflow that other research branches can use. If a future article
about OCR, tokenization, Unicode readiness, or design grammar cannot complete
this checklist, it should remain a research branch rather than a paper
candidate. If it can complete the checklist, it becomes a candidate for the same
evidence-preserving standard used here. The value of this paper is therefore not
only the local result; it is the reusable discipline it imposes on future
research.

\section{Minimum Acceptance Claim}
The minimum acceptance claim is intentionally plain: the paper gives cultural-
heritage computing projects a checked way to keep an artifact-derived result
useful without overstating its source status. That is the claim reviewers should
evaluate first.


\section{Systems Implementation Path}
The invariant can be implemented as a small policy layer around any generative
system that writes, transforms, visualizes, or packages artifact-derived claims.
Each input claim carries a layer label and a maximum public claim. Each transform
declares whether it preserves, weakens, or requires review to strengthen that
label. The output is accepted only when the declared transform is compatible
with the input layer. If the transform tries to strengthen a derived or blocked
claim without a review event, the system rejects the output or rewrites it as a
limitation.

This implementation path is intentionally modest. It does not require a new
large model, a new ontology, or a complete source corpus. It requires that
generation be surrounded by typed evidence checks. In practice, that means a
manuscript writer, visual atlas, product prototype, and benchmark runner can all
share the same rule: no output may exceed the evidence layer it inherited. The
rule is simple enough to test and strict enough to prevent the most damaging
failure seen in this research agenda.

The implementation also creates a useful audit trail. When a generated output is
accepted, the system can record the input layer, transform, output layer, and
review status. When an output is rejected, the system can record the attempted
promotion and the safer rewrite. Over time, this produces a dataset of real
failure modes: which prompts collapse layers, which visual formats hide
uncertainty, which manuscript sections overstate evidence, and which review
instructions actually reduce promotion errors. That dataset is the bridge from a
single invariant to a benchmarkable research agenda.

The audit trail also gives future reviewers something concrete to reject or
improve. If the rejected outputs are trivial, the benchmark is weak. If the safe
rewrites remove useful design capacity, the invariant is too conservative. If
models continue to promote layers despite the policy, the system needs stronger
controls. These are productive failure modes because they turn criticism into
measurable next experiments.

The systems contribution is therefore both a rule and a measurement agenda for
future autonomous research runs.

That agenda is concrete, inspectable, and directly testable.

It can guide the next benchmark iteration immediately.

\section{Benchmark Expansion for Journal Review}
The expanded test is still controlled and synthetic, but it is now large enough
to make the review question concrete. The 24 cases include safe source
statements, safe derived statements, safe weakening, speculative design
hypotheses, blocked public-release claims, and promotions hidden in abstract,
title, contribution, and figure-caption contexts. The comparison includes no
labels, provenance-only checking, citation-only checking, and the full
non-promotion rule. It reports promotion-error detection and baseline false
negatives rather than generic model quality. A rule that blocks every derived
claim would be safe but useless; the current rule preserves safe derived and
weakening cases while rejecting intentional promotions.

The expanded benchmark should use at least four outcome categories. A true
positive is an unsafe promotion that the rule rejects. A true negative is a
safe preservation or weakening that the rule allows. A false positive is a
useful derived or speculative claim that the rule blocks even though the output
keeps its layer. A false negative is a promotion that escapes the rule. The
important metric is not generic model quality. It is severity-weighted
promotion-error recall under a constraint that useful derived design remains
possible. This makes the paper reviewable as a systems result rather than a
philosophical warning.

The benchmark also needs adversarial examples. Some unsafe outputs will not use
obvious phrases such as "source-observed" or "proven." They may imply source
status through table headings, interface labels, graph node colors, or abstract
phrasing. A robust benchmark should therefore include prose, table rows, figure
captions, metadata fields, and product-interface labels. It should ask whether
the invariant catches promotions expressed through structure as well as through
sentences. This is where the method can become more than a local rule: it can
define a reusable evaluation surface for artifact-derived generative systems.

Finally, the benchmark must preserve the human-authority boundary. The rule can
reject a blocked public-release claim, but it cannot decide that a community has
authorized release. It can force a claim into a limitation or review-required
state; it cannot replace review. That separation is the reason the invariant is
useful for culturally sensitive artifacts. It makes generative systems less
likely to overclaim while leaving authority with the people and evidence records
that actually hold it.

\section{Limitations}
The experiment is a controlled case-matrix evaluation, not a live deployment
study. It does not prove that every cultural-heritage AI workflow can be made
safe by label checking, and it does not replace provenance, citation, source
review, or community authority. source transcription review remains required
before any stronger source-layer statement, and rights review remains required
before any public source-image, corpus, or authority claim.

The result is currently a systems paper over claims and
transforms. It is not a public circulation of source images, a glyph-shape grammar,
a Unicode proposal, or a cultural authority decision. It depends on the current
accepted source foundation. If future source review changes the 26 by 8 ledger
or the f/v interpretation, the transform tests must be rerun and the paper must
be revised.

\section{Reproducibility}
The evidence record identifies the manuscript source, count ledger,
article-specific result, transfer audit, rights boundary, and source-review
questions needed to inspect the claim without treating the paper as authorized
for source circulation.

The record includes the expanded layer-safety result, the transfer and
counterexample audit, the non-promotion lemma, the manuscript source, the claim
audit, and the source and rights boundary notes. Reproduction means checking the
24-case matrix, checking the ten transfer cases, verifying that unsafe
promotions are rejected, and confirming that no manuscript statement upgrades
derived or speculative layers into source observations.

\section{Conclusion}
Nwagu Aneke can inspire novel systems only if generated systems preserve the
evidence layers they inherit. This paper contributes a non-promotion invariant
for that purpose. The source-observed foundation remains 26 by 8 equals 208.
The 27/216 layer remains derived by f/v split only. From that foundation,
artifact-derived generative systems can explore applications without converting
design speculation into source evidence.

\bibliographystyle{ACM-Reference-Format}
\bibliography{../references}
\end{document}



